\chapter{数据结构优化 DP}

\section{空间优化}

由于绝大多数情况下 DP 都是在一个 DAG 上转移，可以对此 DAG 分层，每次只保留当前层以及转移需要的层来优化空间复杂度。

\begin{example}[滚动数组]
    若 $f[i][]$ 的转移只需要 $i-k\le j<i$ 的 $j$，则只需要存储这 $k+1$ 层（有时是 $k$ 层）的状态，据此将第一维从 $n$ 优化至 $k$。
\end{example}

如果问题还需要输出方案，那么我们可以通过分治或者分块的方式还原方案。

\begin{example}[滚动数组输出方案]
    P7685 [CEOI 2005] Mobile Service
\end{example}
\begin{solution}[分块]
    按 $B$ 分块，记录第一维在所有 $kB$ 位置的 DP。每次计算一块的方案，首先重算这 $B$ 个位置的 DP 结果，然后倒着算出这一块的方案。据此将第一维的空间从 $n$ 优化至 $O(\sqrt{n})$。
\end{solution}
\begin{solution}[分治]
    $\mathrm{solve}(l,r,S_l,S_r)$ 表示已知一个最优解在 $l$ 时状态为 $S_l$，在 $r$ 时状态为 $S_r$，需要还原这段区间的方案。首先进行一遍 DP 得到这个最优解在 $m=(l+r)/2$ 时刻的状态 $S_m$，然后递归两边即可。
    
    可以将第一维的空间从 $n$ 优化至 $\log n$，但是时间往往需要多一个 $\log n$。
\end{solution}




\section{数据结构优化填表}

本节中，我们考虑最常见的一维 DP 填表转移 $f_i=\bigodot_{j\in [l_i,r_i]} g(j,i)$。对于多维的情况，其思想也是适用的。

在绝大多数情况下，这本质上是一个\textbf{数据结构问题}：我们要查询一个范围中 $g(j,i)$ 的加和，然后计算出 $f_i$，再考虑 $f_i$ 对后续的影响。例如：
\begin{itemize}
    \item $g(j,i)$ 可以分解为 $j$ 的贡献 $a_j$ 与 $i$ 的贡献 $b_i$，此时计算 $f_i$ 相当于查询 $\bigodot_{j\in [l_i,r_i]} a_j$，然后根据 $f_i$ 计算出 $a_i$。
    \item 将 DP 视为关于 $i$ 的扫描线，在 $i\to i+1$ 时维护 $g(\cdot,i)$。
\end{itemize}

\subsection{求和型}

优化 $f_i = \sum_{j\in [l_i,r_i]} g(j,i)$ 类型的 DP。根据 $l_i,r_i$ 和 $g$ 的性质，可以使用不同的数据结构优化。

\begin{example}
    设 $g_{j,i}$ 可以拆成 $O(1)$ 个 $a_ib_j$ 之和，则可以对 $b$ 做前缀和来 $O(1)$ 查询。
\end{example}

\begin{exercise}
    每个 $i$ 向 $[l_i,r_i]$ 连有向边，保证 $i<l_i\le r_i$。求 $1$ 到每个点的路径数量。
\end{exercise}

\begin{example}
    设 $g_{j,i}$ 可以拆成 $O(1)$ 个 $a_ib_jc_{i-j}$ 之和，且 $[l_i,r_i]$ 与 $[1,i)$ 的差别只有 $O(1)$ 项，则可以使用半在线卷积计算，此时 $O(n\log^2 n/\log\log n)$ 是一个常见的复杂度。
\end{example}



\subsection{最值型}

优化 $f_i = \max_{j\in [l_i,r_i]} g(j,i)$ 类型的 DP。根据 $l_i,r_i$ 和 $g$ 的性质，可以使用不同的数据结构优化。由于失去可可减性，这类问题通常比求和型问题更加困难。

\begin{exercise}
    考虑 $g(j,i)=a_j+b_i$ 的形式。
    \begin{itemize}
        \item 使用双端队列在 $O\left(\sum_i |l_i-l_{i-1}|+|r_i-r_{i-1}|\right)$ 复杂度内维护。
        \item 使用 RMQ 来 $O(1)$ 查询 $f_i$。
    \end{itemize}
\end{exercise}

\begin{example}[{P1848 [USACO12OPEN] Bookshelf G}]
\end{example}

\begin{exercise}
    每个 $i$ 向 $[l_i,r_i]$ 连带权有向边，保证 $i<l_i\le r_i$。求 $1$ 到每个点的最长路长度。
\end{exercise}

\section{整体 DP}

有时我们需要处理一些多维的 DP，其中某一维上的转移基本都是原地且相同的。对于这种问题，我们可以考虑直接使用数据结构的区间/全局操作来维护这一维。

\begin{example}[{P9400 「DBOI」Round 1 三班不一般}]
    支持以下转移：
    \begin{itemize}
        \item 对 $0<j<A$，$f_{i,j}=a_if_{i-1,j-1}+b_i$
        \item $f_{i,0}=a_i\sum_j f_{i-1,j}$
    \end{itemize}
\end{example}

\begin{example}
    
\end{example}

\section{矩阵乘法}

有时，DP 转移可以视为一个矩阵乘法。此时可以使用矩阵快速幂来询问远处点值，或者用数据结构维护矩阵乘法乘积来支持修改或区间询问。
\begin{remark}
    对于询问远处点值的情况，矩阵快速幂通常可以被一个线性递推替代。
\end{remark}

\begin{example}[{P6772 [NOI2020] 美食家}]
    
\end{example}

\begin{example}[{P4719 【模板】动态 DP}]
    
\end{example}

\section{凸包}

\subsection{凸集、凸包、凸函数}

\begin{definition}[凸组合]
    对于 $\mathbb R^d$ 中的若干个点 $\mathbf x_i$ 的线性组合 $\sum_i \theta_i \mathbf x_i$，若系数 $\theta$ 满足 $\sum_i \theta_i=1$ 且 $\theta_i\ge 0$，则称此线性组合为一个\textbf{凸组合}。
\end{definition}

\begin{definition}[凸集]
    对于点集 $S\subseteq \mathbb R^d$，如果其中任何两个点 $\mathbf x_1,\mathbf x_2$ 的所有凸组合 $\theta \mathbf x_1+(1-\theta)\mathbf x_2$ 都在 $S$ 中，则称 $S$ 是凸的。换句话说，$S$ 是凸的当且仅当连接其中任意两点的线段都被包含在 $S$ 中。
\end{definition}

\begin{property}
    \begin{itemize}
        \item 凸集的交是凸集
        \item 凸集的 Minkowski 和 $\{\mathbf x+\mathbf y:\mathbf x\in S_1,\mathbf y\in S_2\}$ 是凸集
        \item 凸集的直和 $\{(\mathbf x,\mathbf y):\mathbf x\in S_1,\mathbf y\in S_2\}$ 是凸集
    \end{itemize}
\end{property}


\begin{definition}[凸包]
    对于点集 $S\subseteq \mathbb R^d$，其\textbf{凸包}（也称\textbf{凸壳}）定义为包含 $S$ 的最小凸集，即所有包含 $S$ 的凸集的交集。
\end{definition}

\begin{theorem}
    设 $S\subseteq \mathbb R^d$。以下表述等价：
    \begin{enumerate}
        \item $C$ 是 $S$ 的凸包
        \item $C$ 是 $S$ 中元素的所有有限凸组合构成的集合，即
            $$
            \left\{\sum_{i=1}^k \theta_i \mathbf x_i:\theta_i\ge 0,\sum_i \theta_i=1,\mathbf x_i\in S\right\}.
            $$
        \item $C$ 是 $S$ 中所有 $d+1$ 个点的凸组合构成的集合
    \end{enumerate}

    若 $S$ 有限，则以下表述等价：
    \begin{enumerate}
        \item $C$ 是 $S$ 的凸包
        \item $C$ 是包含 $S$ 的最小凸多面体
        \item $C$ 是所有包含 $S$ 的半超平面的交
    \end{enumerate}
\end{theorem}

\begin{fact}
    点 $\mathbf{p}$ 是 $S$ 的凸包顶点当且仅当它不是 $S-\{\mathbf p\}$ 中点的凸组合。因此，对于有限点集 $S$，凸包的顶点数量不超过 $|S|$。
\end{fact}

\begin{lemma}
    对于有限点集 $S\subseteq \mathbb R^d$，令 $C$ 为其凸包。对于任何 $\mathbf w\in \mathbb R^d$，都有
    $$
    \max_{\mathbf x\in S}\mathbf w\cdot \mathbf x=\max_{\mathbf x\in \partial C}\mathbf w\cdot \mathbf x.
    $$
\end{lemma}
\begin{proof}
    由于 $S$ 有限，$C$ 是闭集，故最大值可以在某个 $x\in \mathbf S$ 中取到。对于 $\mathbf x\in \mathrm{int}(C)$，必然可以将其向某一方向移动 $\epsilon$ 距离，使得移动后仍然处于 $C$ 中，且 $\mathbf w\cdot \mathbf x'$ 增大。故最大值只能在 $\partial C$ 上取到。
\end{proof}

在绝大多数情况下，我们只需要关注二维情况。此时，凸包是一个凸多边形。
\begin{definition}[上凸壳和下凸壳]
    对于点集 $S\subseteq \mathbb R^2$，设 $C$ 为其凸包。考虑 $C$ 中 $x$ 坐标最小和最大的点 $L,R$，顺时针沿着凸包边界从 $L$ 走到 $R$ 得到的路径称为\textbf{上凸壳}，从 $R$ 走到 $L$ 得到的路径称为\textbf{下凸壳}。
\end{definition}
\begin{corollary}
    考虑点集 $S\subseteq \mathbb R^2$。对于给定的 $k$，$\max_{(x,y)\in S}kx+y$ 在上凸壳中取到，$\min_{(x,y)\in S}kx+y$ 在下凸壳中取到。$(x_0,y_0)\in S$ 处可以取得最值当且仅当直线 $y-y_0=k(x-x_0)$ 是凸包的切线。
\end{corollary}
\begin{exercise}
    对 $\min_{(x,y)\in S}kx+y$ 给出类似的结论。
\end{exercise}

\begin{definition}[凸函数]
    若定义在凸集 $C$ 上的函数 $f:C\to \mathbb R$ 满足对任何 $\mathbf x_1,\mathbf x_2\in C$ 和 $\theta\in (0,1)$，都有
    $$
    f(\theta \mathbf x_1+(1-\theta)\mathbf x_2)\le \theta f(\mathbf x_1)+(1-\theta)f(\mathbf x_2)
    $$
    成立，则称 $f$ 是\textbf{凸(convex)}的。若不等号严格成立，则称 $f$ 是\textbf{严格凸(strictly convex)}的。若不等式反向成立，则称 $f$ 是\textbf{凹(concave)}的。
\end{definition}
\begin{remark}
    在国内的一些教材中，凸和凹的定义是反过来的。为了避免混淆，你也可以使用“上凸”和“下凸”来代替“凸”和“凹”。在本讲义中，我们遵循国际惯例，使用“凸”来表示“下凸”，使用“凹”来表示“上凸”。
\end{remark}
\begin{remark}
    在这个定义中，我们并不要求函数可导。
\end{remark}



\begin{property}
    设 $f:C\to \mathbb R$ 是一个凸函数。有以下性质成立：
    \begin{itemize}
        \item 切线在图像下方：假设 $f$ 在 $\mathrm{int}(C)$ 上可导。对任何 $\mathbf x_1,\mathbf x_2\in \mathrm{int}(C)$，都有 $f(\mathbf x_2)\ge f(\mathbf x_1)+\nabla f(\mathbf x_1)\cdot (\mathbf x_2-\mathbf x_1)$。反过来，满足该性质的函数是一个凸函数；若进一步满足不等号严格成立，则函数是严格凸的。
        \begin{itemize}
            \item 对于一维情况，就是 $f(x_2)\ge f(x_1)+f'(x_1)(x_2-x_1)$。
        \end{itemize}
        \item 假设 $f$ 在 $\mathrm{int}(C)$ 上二阶可导。对任何 $\mathbf x$，$\nabla^2 f(\mathbf x)$ 是半正定的。反过来，满足该性质的函数是一个凸函数；若进一步满足不等号严格成立，则函数是严格凸的。
        \begin{itemize}
            \item 对于一维情况，即 $f''(x)\ge 0$。这也可以理解为斜率不降：对任何 $x<y<z$，都有 $\frac{f(y)-f(x)}{y-x}\le \frac{f(z)-f(x)}{z-x}\le \frac{f(z)-f(y)}{z-y}$。
        \end{itemize}
        \item $f$ 的\textbf{上图像(epigraph)} $\mathrm{epi}(f)=\{(\mathbf x,y):y\ge f(\mathbf x)\}$ 是凸集。反过来，满足该性质的函数是凸的。
        \item Jensen 不等式：对于任意 $\mathbf x_1,\ldots,\mathbf x_n\in C$ 的凸组合 $\overline{\mathbf x}=\sum_{i=1}^n \theta_i \mathbf x_i$，都有
        $$
        f\left(\sum_{i=1}^n \theta_i \mathbf x_i\right)\le \sum_{i=1}^n \theta_i f(\mathbf x_i).
        $$
    \end{itemize}
\end{property}

\begin{property}
    \begin{itemize}
        \item 凸函数的非负线性组合是凸的
        \item 凸函数的 $\max$ 是凸的
        \item 凸函数复合一次函数是凸的
        \item 设 $f,g:\mathbb R\to \mathbb R$ 是凸的，且 $g$ 在 $f(\mathbb R)$ 上不降，则 $g\circ f$ 是凸的。
    \end{itemize}
\end{property}

\begin{example}
    \begin{itemize}
        \item 一次函数 $ax+b$ 同时是凸的和凹的
        \item 开口向上的二次函数 $a(x-b)^2+c$ 是凸的
        \item 绝对值函数 $|x-a|$ 是凸的
    \end{itemize}
\end{example}

以上讨论的都是连续情况下的定义，离散情况下则并没有一个统一的定义。我们使用一个相对常用的定义：
\begin{definition}[离散凸函数]
    对于函数 $f:\{l,l+1,\ldots,r\}\to \mathbb R$，若其二阶差分 
    $$
    \Delta^2 f(x)=f(x+2)-2f(x+1)+f(x)
    $$
    非负，则称 $f$ 是凸的。若 $\Delta^2 f(x)>0$，则称 $f$ 是严格凸的。
\end{definition}

\begin{exercise}
    对于定义域包含 $\{l,\ldots,r\}$ 的连续凸函数 $f$，证明：$f$ 限制到 $\{l,\ldots,r\}$ 上得到的函数是离散凸的。
\end{exercise}

\begin{property}
    设 $f:\{l,\ldots,r\}\to \mathbb R$ 的凸的。以下命题等价：
    \begin{itemize}
        \item $f$ 是凸的
        \item $f$ 的斜率不降：对任何 $l\le x<y<z\le r$，都有 $\frac{f(y)-f(x)}{y-x}\le \frac{f(z)-f(x)}{z-x}\le \frac{f(z)-f(y)}{z-y}$。
        \item $f$ 的一阶差分 $\Delta f(x)=f(x+1)-f(x)$ 不降
        \item 切线在函数下方：对任何 $l\le x,y\le r,\Delta f(x-1)\le k\le \Delta f(x)$，都有 $f(y)\ge f(x)+k(y-x)$。
        \item 延拓到连续函数后凸：对每个 $i,i+1$，用一次函数连接 $(i,f(i))$ 和 $(i+1,f(i+1))$ 得到的 $[l,r]$ 上的连续函数 $\tilde f$ 是凸的。
    \end{itemize}
\end{property}
利用最后一个命题，我们可以在连续凸和离散凸之间互相转化。

\subsection{二维凸包}

本节中，我们介绍二维凸包的算法和相关应用。

\begin{lemma}
    考虑点集 $S$ 的凸壳上的所有顶点，一个点 $\mathbf x$ 顺/逆时针方向的下一个点 $\mathbf y$ 是使得 $\mathbf y-\mathbf x$ 的极角最大/小的 $\mathbf{y}\in S$。
\end{lemma}
\begin{remark}
    这里大/小的定义为：在所有可能的跨度不超过 $\pi$ 的极角中，逆时针的最后一个/第一个。
\end{remark}

\begin{descbox}{算法}[Andrew 算法]
    Andrew 算法从左往右分别构建出上、下凸壳。将所有点按照 $(x,y)$ pair 排序。从左到右扫描加入点，维护所有线段的上、下包络，每次弹出末尾的一段点。排序后线性。
\end{descbox}
\begin{remark}
    常常只需要上/下凸壳中的一个。
\end{remark}
\begin{remark}
    一般来说，不需要保留线段的中间点。在需要保留所有边界点时，可以将弹栈的条件从严格不等号改为非严格不等号。
\end{remark}

\begin{descbox}{算法}[Graham 算法]
    Graham 算法按照极角序构建凸包。首先选择 $(x,y)$ 按 pair 比较时最小的一个，该点必然在凸包上，且所有其他点都在以其为原点的 $(-\frac{\pi}{2},\frac{\pi}{2}]$ 角度范围内。按照角度排序后逆时针构建凸包即可。
\end{descbox}

\begin{example}
    维护平面上的点集 $S$，支持以下操作：
    \begin{itemize}
        \item 插入一个点 $(x_i,y_i)$。保证 $x$ 不降。
        \item 给定 $k$，询问 $\max_{(x,y)\in S}kx+y$。
    \end{itemize}
\end{example}
\begin{exercise}
    假设询问的 $k$ 不增/不降，说明取到最值的点只会在凸包上单调移动。据此设计一个 $O(n)$ 的算法。
\end{exercise}
\begin{exercise}
    优化 $f_i=\min_{j<i}k_i\cdot a_j+b_j$ 类型的 DP。当 $k_i$ 不降时，这被称为\textbf{斜率优化}。

    P3195 [HNOI2008] 玩具装箱
\end{exercise}

\begin{example}
    在上一题中去掉 $x$ 不降的限制。
\end{example}
\begin{solution}
    平衡树维护凸包/李超树/离线对时间分治。李超树比较好写且快。
\end{solution}
\begin{remark}
    实际上还可以支持删除，这就是完全动态凸包问题。
\end{remark}

\begin{example}[{P3309 [SDOI2014] 向量集}]
    维护序列 $(x_i,y_i)$，支持以下操作：
    \begin{itemize}
        \item 在末尾插入一个点 $(x_i,y_i)$。
        \item 给定 $l,r,a,b$，询问 $\max_{i\in [l,r]}ax_i+by_i$。
    \end{itemize}
\end{example}
\begin{exercise}
    额外支持末尾删除
\end{exercise}
\begin{exercise}
    优化 $f_i=\min_{j\in [l_i,r_i]}k_i\cdot a_j+b_j$ 类型的 DP。
\end{exercise}

在很多情况下（如背包问题），我们会使用集合的卷积 $H_i=\cup_j (F_j+G_{i-j})$，其中 $S+T=\{x+y:x\in S,y\in T\}$。注意这相当于对 $F=\{(i,x):x\in F_i\}$ 和 $G=\{(i,y):y\in G_i\}$ 做 Minkowski 和得到 $H=\{(i,z):z\in H_i\}$。记 $C(S)$ 为点集 $S$ 的凸包，则有以下性质成立：
\begin{property}
    \begin{itemize}
        \item $C(S\cup T)=C(C(S)\cup C(T))$
        \item $C(S+T)=C(S)+C(T)$
    \end{itemize}
\end{property}

于是，对点集进行运算后再求凸包可以变为对点集的凸包进行运算。并集的凸包容易计算，而 Minkowski 和的凸包可以通过以下算法计算。
\begin{descbox}{算法}[凸包的 Minkowski 和]
    只需计算两个凸包顶点集合 $C_1,C_2$ 的 Minkowski 和。令 $\mathbf p_1,\mathbf p_2$ 分别为 $C_1,C_2$ 中按 $(x,y)$ 比较时最小的点，则 $\mathbf p=\mathbf p_1+\mathbf p_2$ 必然为 $S_1+S_2$ 中 $(x,y)$ 最小的点。

    从 $\mathbf p$ 出发逆时针逐渐构建新凸包的顶点，设目前的点为 $\mathbf x_1+\mathbf x_2$，下一个顶点 $\mathbf x'_1+\mathbf x'_2$ 必然是关于当前顶点极角最小的一个，而这只有两种可能：
    \begin{itemize}
        \item $\mathbf x'_1$ 是 $\mathbf x_1$ 在 $C_1$ 中的下一个点，$\mathbf x'_2=\mathbf x_2$
        \item $\mathbf x'_1=\mathbf x_1$，$\mathbf x'_2$ 是 $\mathbf x_2$ 在 $C_2$ 中的下一个点
    \end{itemize}
    在其中选择较小的一个即可。复杂度 $O(|C_1|+|C_2|)$。
\end{descbox}
\begin{remark}
    可以发现，这相当于把两个凸包的边“按极角归并”。
\end{remark}


考虑函数的 $(\min,+)$ 卷积与逐点 $\min$ 的上图像，有以下事实成立： 
\begin{fact}
    考虑定义在整数上的函数 $f,g$。
    \begin{itemize}
        \item 考虑其 $(\min,+)$ 卷积 $h(i)=\min_{j}f(j)+g(i-j)$。则 $\mathrm{epi}(h)=\mathrm{epi}(f)+\mathrm{epi}(g)$。
        \item 考虑其逐点最小值 $h(i)=\min(f(i),g(i))$。则 $\mathrm{epi}(h)=\mathrm{epi}(f)\cup \mathrm{epi}(g)$。
    \end{itemize}
\end{fact}
考虑由若干函数进行一系列逐点 $\min$ 和 $(\min,+)$ 卷积得到的函数 $f$。有时，我们可以说明 $f$ 是凸的；有时，我们只需要 $\mathrm{epi}(f)$ 的凸包（如只需要计算 $F(x)=\min_i f(i)+ix$）。此时，我们可以利用前述的点集凸包的运算来进行优化。首先，我们有以下事实：
\begin{fact}
    $\mathrm{epi}(f)$ 的凸包边界等于 $\{(i,f(i))\}$ 的下凸壳。
\end{fact}
于是，我们只需要维护过程中的所有下凸壳，对下凸壳求并集或 Minkowski 和即可。

\begin{example}
    在线性时间内计算凸函数的 $(\min,+)$ 卷积。
\end{example}
\begin{exercise}
    在线性时间内计算凸函数的 $(\max,+)$ 卷积。
    \begin{hint}
        固定 $k$，$H_k(i)=f_i+g_{k-i}$ 是凸的。
    \end{hint}
\end{exercise}

\begin{example}[{P5073 [Ynoi Easy Round 2015] 世上最幸福的女孩}]
    全局加，询问区间最大子段和。离线。
\end{example}



% \begin{table}[h]
% \centering
% \begin{tabular}{@{}llll@{}}
% \toprule
% 数据结构 & 预处理 & 查询 & 适用范围 \\ \midrule
% 前缀和 & $O(n)$ & $O(1)$ & 具有逆元 (群) \\
% ST 表 & $O(n \log n)$ & $O(1)$ & 可重复贡献 (半格) \\
% 线段树 & $O(n)$ & $O(\log n)$ & 一般半群 \\
% Disjoint ST & $O(n \log n)$ & $O(1)$ & 一般半群 \\
% \bottomrule
% \end{tabular}
% \end{table}